\chapter{Atualidades e Inteligência Artificial}
\label{cap:atualidades-teoria}

Este capítulo cobre as duas partes do bloco de Atualidades: o viés
socioambiental clássico de prova de conhecimentos gerais, e o bloco novo de
Inteligência Artificial (que entrou no edital da Dataprev só em 2026). Como
é o capítulo mais volátil do livro --- números, status de projeto de lei e
métricas de mercado mudam --- o objetivo aqui é fixar o \textbf{mecanismo}
de cada conceito (por que ESG mede o que mede, por que viés e variância se
comportam de forma oposta, por que um LLM alucina) para que o entendimento
sobreviva mesmo quando o número específico for atualizado antes da prova.

\section{Atualidades: o viés socioambiental}

\begin{conceito}[ESG --- avaliando o que o balanço financeiro não mostra]
\textbf{ESG} (\textit{Environmental, Social and Governance} ---
\textbf{A}mbiental, \textbf{S}ocial e \textbf{G}overnança) é o conjunto de
critérios usado para avaliar o desempenho \textbf{não financeiro} de uma
organização --- o que o balanço contábil tradicional não captura, mas que
afeta o risco e a reputação de longo prazo do negócio:

\begin{itemize}
\item \textbf{Ambiental (E)}: emissões de carbono, consumo de energia e
água, gestão de resíduos, biodiversidade, exposição a mudanças climáticas.
\item \textbf{Social (S)}: condições de trabalho, diversidade e inclusão,
direitos humanos, relação com comunidades afetadas, práticas na cadeia de
fornecedores.
\item \textbf{Governança (G)}: composição e independência do conselho de
administração, ética corporativa, transparência, combate à corrupção,
prestação de contas aos acionistas e à sociedade.
\end{itemize}

O motivo de existir como critério \textbf{formal} de investimento e
relatório corporativo: investidores e reguladores perceberam que riscos
não financeiros (um desastre ambiental, um escândalo de corrupção) também
destroem valor financeiro --- então passaram a exigir que esse risco seja
\textbf{medido e divulgado}, não só descoberto depois que já aconteceu.
\end{conceito}

\begin{cuidado}[O “G” não é “Gestão”, “Global” nem “Growth”]
“Governança” é especificamente sobre \textbf{como a organização é dirigida
e controlada} (estrutura de poder, prestação de contas) --- não é
sinônimo genérico de “boa gestão” nem está relacionado a expansão
(“Growth”) ou alcance internacional (“Global”). Da mesma forma, o “S” é
\textbf{Social}, não “Sustentável” --- sustentabilidade é o conceito
guarda-chuva que o ESG inteiro serve, não uma das três letras
especificamente.
\end{cuidado}

\begin{regrapratica}[Como filtrar item de atualidades socioambientais]
Esse tipo de questão costuma vir ancorada num \textbf{texto-base}, no
formato de julgar afirmativas (I, II, III). Dois filtros rápidos, antes de
reler o texto com calma: (1) item com termo \textbf{absoluto}
(“invariavelmente”, “sempre”, “nunca”, “exclusivamente”) é quase sempre
\textbf{falso} --- a realidade socioambiental raramente é tão categórica;
(2) item com dado \textbf{numérico} precisa ser \textbf{conferido contra o
texto-base}, nunca aceito de memória --- é comum a alternativa inflar ou
distorcer um número que o texto deu de forma mais modesta.
\end{regrapratica}

\section{Fundamentos de Inteligência Artificial}

\subsection{IA, Machine Learning e Deep Learning: uma relação de contêiner}

\begin{conceito}[Três círculos concêntricos, não três coisas separadas]
\begin{itemize}
\item \textbf{Inteligência Artificial (IA)}: o campo mais amplo ---
qualquer sistema capaz de executar tarefas que, se feitas por um humano,
exigiriam inteligência (raciocinar, perceber, decidir).
\item \textbf{Machine Learning (ML)}: um \textbf{subcampo} da IA em que o
sistema \textbf{aprende padrões a partir de dados} e melhora seu
desempenho com a experiência, em vez de seguir regras explicitamente
programadas por um humano para cada situação.
\item \textbf{Deep Learning}: um \textbf{subcampo do ML} que usa
\textbf{redes neurais com muitas camadas} (daí “profundo”,
\textit{deep}) para aprender representações cada vez mais abstratas dos
dados, camada após camada.
\end{itemize}
A relação é de \textbf{contenção}: IA $\supset$ ML $\supset$ Deep
Learning --- todo Deep Learning é ML, todo ML é IA, mas a recíproca não
vale (existe IA que não usa aprendizado de máquina, como sistemas
baseados puramente em regras lógicas escritas à mão).
\end{conceito}

\subsection{Tipos de aprendizado}

\begin{conceito}[O que diferencia os três --- o que o sistema recebe como sinal de aprendizado]
\begin{center}
\begin{tabular}{@{}p{2.6cm}p{2.8cm}p{2.6cm}p{3cm}@{}}
\toprule
\textbf{Tipo} & \textbf{Dados de treino} & \textbf{Objetivo} & \textbf{Exemplos} \\
\midrule
Supervisionado & \textbf{rotulados} (resposta certa conhecida) & prever o
rótulo de casos novos & classificação, regressão \\
Não supervisionado & \textbf{sem rótulo} & achar estrutura escondida nos
dados & \textit{clustering} (K-Means), associação, PCA \\
Por reforço & recompensa/punição a cada ação & aprender uma
\textbf{política} de ação ótima & agentes, jogos, robótica \\
\bottomrule
\end{tabular}
\end{center}

O critério que separa os três é \textbf{que tipo de sinal} orienta o
aprendizado: no supervisionado, o sinal é “aqui está a resposta certa
para este exemplo”; no não supervisionado, não há resposta certa dada,
só os dados brutos; no aprendizado por reforço, o sinal é indireto ---
“essa sequência de ações rendeu recompensa X”, sem dizer qual ação
específica foi a certa.
\end{conceito}

\begin{cuidado}[“Aprende com os dados e melhora ao longo do tempo” não é lógica booleana]
Um enunciado que descreve um sistema “que aprende com a experiência e
melhora seu desempenho ao longo do tempo” está descrevendo \textbf{redes
neurais/aprendizado de máquina} --- não lógica booleana clássica nem
programação linear determinística, que executam sempre a mesma regra fixa
e não mudam de comportamento sozinhas com o uso.
\end{cuidado}

\section{Viés e variância: o trade-off central do aprendizado supervisionado}

\begin{conceito}[Dois jeitos opostos de um modelo errar]
Todo modelo supervisionado comete um dos dois tipos de erro sistemático
(ou uma combinação dos dois), e eles têm causas \textbf{opostas}:

\begin{center}
\begin{tabular}{@{}p{2.6cm}p{4.4cm}p{4.4cm}@{}}
\toprule
& \textbf{Alto viés (\textit{bias})} & \textbf{Alta variância (\textit{variance})} \\
\midrule
Causa & modelo \textbf{simples demais} para captar o padrão real dos
dados & modelo \textbf{complexo demais}, ajusta-se a ruído específico do
treino \\
Sintoma & erra \textbf{tanto no treino quanto no teste} --- nem decorou o
treino direito & acerta muito bem o treino, erra feio em dados novos ---
“decorou” em vez de generalizar \\
Nome do fenômeno & \textbf{underfitting} & \textbf{overfitting} \\
Como combater & modelo mais expressivo, mais atributos relevantes &
regularização (L1/L2), validação cruzada, mais dados de treino \\
\bottomrule
\end{tabular}
\end{center}

O motivo de ser um \textbf{trade-off} (e não duas falhas independentes que
dá para zerar ao mesmo tempo): reduzir a complexidade do modelo para
combater a variância tende a aumentar o viés (o modelo fica simples
demais de novo); aumentar a complexidade para combater o viés tende a
aumentar a variância (o modelo passa a se ajustar a detalhes específicos
do treino). O objetivo prático não é zerar os dois, é achar o ponto de
\textbf{equilíbrio} que minimiza o erro total em dados novos.
\end{conceito}

\begin{cuidado}[Não inverta viés com variância]
“Viés alto” \textbf{não} é o modelo que “decorou” o treino --- isso é
\textbf{variância alta} (overfitting). O jeito de fixar sem trocar: viés
alto = modelo \textbf{“burro”} (simples demais até para o próprio treino);
variância alta = modelo que \textbf{“decorou”} (bom no treino, ruim fora
dele). E regularização combate especificamente a \textbf{variância/
overfitting}, não o viés --- regularizar um modelo que já sofre de viés
alto (underfitting) só pioraria a situação, tornando-o ainda mais
restrito.
\end{cuidado}

\section{Métricas de avaliação}

\begin{conceito}[Classificação: quatro métricas a partir da matriz de confusão]
\begin{itemize}
\item \textbf{Acurácia} = acertos totais / total de casos. Parece a
métrica mais natural, mas é \textbf{enganosa em bases desbalanceadas}.
\item \textbf{Precisão}: dos casos que o modelo \textbf{apontou como
positivos}, quantos realmente eram --- importa quando o custo de um
\textbf{falso positivo} é alto (ex.: barrar um e-mail legítimo como se
fosse spam).
\item \textbf{Recall} (revocação/sensibilidade): dos casos que
\textbf{realmente eram positivos}, quantos o modelo conseguiu capturar ---
importa quando o custo de um \textbf{falso negativo} é alto (ex.: deixar
passar uma fraude).
\item \textbf{F1-score}: a \textbf{média harmônica} entre precisão e
recall --- harmônica, e não aritmética, justamente porque a média
harmônica \textbf{pune mais} o desequilíbrio entre as duas (um modelo com
precisão 100\% e recall 1\% teria média aritmética de 50,5\%, uma nota
enganosamente alta; a harmônica cairia para perto de 2\%, refletindo
melhor que o modelo é inútil na prática).
\end{itemize}
\end{conceito}

\begin{exemplo}[Por que acurácia engana em base desbalanceada]
Um conjunto de e-mails tem 99\% legítimos e 1\% spam. Um modelo “burro”
que \textbf{sempre} responde “legítimo”, sem olhar o conteúdo, acerta
\textbf{99\% de acurácia} --- um número que parece excelente --- mas
detecta \textbf{zero} spam (recall igual a zero para a classe que
realmente importa detectar). É exatamente o cenário que a prova costuma
montar (spam, fraude, doença rara) para testar se o candidato confia
demais na acurácia isolada.
\end{exemplo}

\begin{conceito}[Regressão: MAE, RMSE e R²]
Para prever um \textbf{número} (não uma classe), as métricas de
classificação (acurácia, precisão, F1) não se aplicam --- usam-se
\textbf{MAE} (erro médio absoluto, na mesma unidade do valor previsto),
\textbf{RMSE} (raiz do erro quadrático médio --- penaliza mais erros
grandes, por elevar ao quadrado antes de tirar a média) e \textbf{R²}
(quanto da variação dos dados o modelo consegue explicar, de 0 a 1). Usar
uma métrica de classificação num problema de regressão (ou vice-versa) é
um distrator comum --- confirme sempre se a saída do modelo é uma
\textbf{categoria} ou um \textbf{número} antes de escolher a métrica.
\end{conceito}

\section{Modelos generativos e LLMs}

\begin{conceito}[Modelo generativo: aprende a distribuição para criar exemplos novos]
Um \textbf{modelo generativo} aprende a \textbf{distribuição estatística}
dos dados de treino e usa esse aprendizado para \textbf{gerar exemplos
novos}, plausíveis dentro daquela distribuição --- texto, imagem, áudio ---
em vez de apenas classificar ou prever um rótulo sobre um exemplo já
existente (o que seria um modelo \textbf{discriminativo}).
\end{conceito}

\begin{conceito}[LLM: linguagem em escala, via Transformer e atenção]
Um \textbf{LLM} (\textit{Large Language Model}) é um modelo de linguagem
treinado sobre um volume massivo de texto, tipicamente com a arquitetura
\textbf{Transformer} e seu \textbf{mecanismo de atenção} --- que permite ao
modelo, ao processar cada palavra, “prestar atenção” dinamicamente em
\textbf{quais outras palavras do contexto} são mais relevantes para aquele
ponto, em vez de processar a sequência de forma estritamente linear e
igualmente ponderada. É esse mecanismo que permite ao modelo capturar
dependências de longo alcance no texto (uma referência no início do
parágrafo que só faz sentido lida junto com uma palavra no final).
\end{conceito}

\begin{conceito}[RAG: buscar antes de gerar]
\textbf{RAG} (\textit{Retrieval-Augmented Generation}) combina o LLM com
uma etapa de \textbf{busca} (\textit{retrieval}) numa base de conhecimento
externa \textbf{antes} de gerar a resposta --- em vez de responder só com
o que aprendeu durante o treinamento (que tem uma data de corte e pode
estar desatualizado ou não cobrir um domínio específico), o modelo recebe
trechos relevantes recuperados da base e os usa como contexto adicional
para fundamentar a resposta. É a técnica típica para tornar respostas mais
\textbf{atuais} e \textbf{rastreáveis} a uma fonte.
\end{conceito}

\begin{conceito}[Alucinação: resposta fluente, mas não fundamentada]
\textbf{Alucinação} é quando o modelo gera uma resposta \textbf{plausível
e bem escrita}, mas \textbf{factualmente incorreta} ou inventada --- o
modelo não “sabe que não sabe”: ele continua produzindo texto
estatisticamente coerente mesmo quando não tem base real para a afirmação.
É a razão pela qual saída de LLM em contexto que afeta decisão real (um
atendimento ao cidadão, uma análise jurídica) exige \textbf{validação},
não confiança cega.
\end{conceito}

\section{Ética, governança e privacidade em IA}

\begin{conceito}[Viés algorítmico: o modelo herda o viés dos dados, não o inventa]
\textbf{Viés algorítmico} ocorre quando os \textbf{dados de treino}
refletem desigualdades ou padrões históricos discriminatórios, e o modelo
\textbf{aprende e reproduz} esse padrão em suas decisões --- o algoritmo
não é “preconceituoso” por conta própria, ele generaliza exatamente o que
os dados de treino mostraram, inclusive quando esses dados carregam
discriminação histórica.
\end{conceito}

\begin{cuidado}[Não usar o atributo sensível diretamente não elimina o viés]
Um modelo que não usa “raça” ou “gênero” como atributo de entrada
\textbf{não} está automaticamente livre de viés --- outras variáveis
(CEP, escola frequentada, histórico de consumo) podem funcionar como
\textbf{proxy}, recriando indiretamente a mesma discriminação através de
uma correlação estatística. A alternativa que descreve isso como “o
modelo é neutro porque não usa o atributo sensível” é o distrator clássico
de “neutralidade tecnológica” --- neutralidade de dado não é garantida
só por omitir uma coluna.
\end{cuidado}

\begin{conceito}[Explicabilidade x acurácia: são eixos independentes]
\textbf{Explicabilidade} é a capacidade de o modelo (ou de quem o opera)
\textbf{justificar} uma decisão de forma compreensível a humanos.
\textbf{Acurácia} é o quanto o modelo acerta. Os dois não andam
necessariamente juntos: um modelo pode ser \textbf{muito acurado e pouco
explicável} (as redes neurais profundas são o exemplo clássico de
“caixa-preta” --- acertam muito, mas é difícil apontar exatamente
\emph{por quê} decidiram assim). Essa tensão é o que justifica exigir
explicabilidade como requisito \textbf{separado} de desempenho em
contextos de decisão sensível (crédito, benefício social, Justiça).
\end{conceito}

\begin{conceito}[LGPD, art.\ 20 --- direito à revisão de decisão automatizada]
A LGPD garante ao titular de dados o direito de \textbf{solicitar revisão}
de decisões tomadas \textbf{unicamente} com base em tratamento
automatizado que afetem seus interesses --- incluindo decisões que
definam perfil pessoal, profissional, de consumo ou de crédito. É a base
legal que ancora cenários de prova envolvendo score de crédito, triagem de
currículo automatizada ou concessão de benefício público por algoritmo:
sistema que decide \textbf{sozinho}, sem qualquer possibilidade de revisão
humana, contraria essa garantia.
\end{conceito}

\section{Regulação da IA: em vigor x ainda em tramitação}

\begin{conceito}[A distinção de estado que a prova mais explora]
Duas normas de referência sobre IA existem hoje em estágios
\textbf{completamente diferentes}, e confundir os dois estágios é o erro
mais fácil de cometer por quem só acompanha manchete de notícia:

\begin{itemize}
\item \textbf{União Europeia --- Regulamento (UE) 2024/1689 (“AI Act”)}:
\textbf{já está em vigor} (desde 01/08/2024), com aplicação
\textbf{escalonada} no tempo (as proibições entraram primeiro; as
obrigações de sistemas de alto risco entram depois). Por ser um
\textbf{regulamento} europeu, aplica-se diretamente nos Estados-membros,
sem precisar de uma lei nacional para “ativá-lo”. Estrutura-se por
\textbf{nível de risco}: risco \textbf{inaceitável} (práticas proibidas,
como pontuação social do cidadão pelo Estado), \textbf{alto risco}
(obrigações pesadas de conformidade e supervisão humana), \textbf{risco
limitado} (dever de transparência --- avisar que se está interagindo com
uma IA) e \textbf{risco mínimo} (sem obrigação especial).
\item \textbf{Brasil --- PL 2338/2023}: \textbf{ainda é projeto de lei}
--- não é lei em vigor. Segue o mesmo desenho baseado em \textbf{nível de
risco} do modelo europeu, prevê \textbf{direitos} das pessoas afetadas
(informação prévia, explicação da decisão, contestação e revisão humana),
cria um sistema nacional de governança de IA sob coordenação da ANPD, e
prevê multa por infração.
\end{itemize}

O mecanismo comum aos dois textos: quanto maior o \textbf{risco} que um
sistema de IA representa para direitos das pessoas, maior a
\textbf{obrigação} de conformidade imposta a quem o opera --- é uma
regulação \textbf{proporcional ao risco}, não uma proibição genérica de
IA.
\end{conceito}

\begin{cuidado}[As trocas de estado que passam despercebidas]
\begin{itemize}
\item Dizer que o \textbf{Brasil já tem} uma lei específica de IA em
vigor é falso --- o que está em vigor no Brasil e toca IA é a
\textbf{LGPD} (por tabela, via o art.\ 20 sobre decisão automatizada), não
uma lei de IA dedicada.
\item Dizer que o \textbf{AI Act ainda é uma proposta} em discussão é
falso --- é regulamento \textbf{vigente} desde agosto de 2024.
\item Atribuir ao PL 2338 uma abordagem “baseada em princípios, sem
classificação de risco” é falso --- ele é, como o europeu,
\textbf{baseado em nível de risco}.
\end{itemize}
Como este é o assunto mais volátil do livro, \textbf{confira o estágio
atual do PL 2338} próximo à data da prova --- se ele for sancionado como
lei antes disso, deixa de ser “projeto” e a resposta correta muda de
categoria.
\end{cuidado}

\section{Reconhecendo o padrão em cenários de IA aplicada}

\begin{conceito}[Quatro padrões de cenário que a prova reaproveita]
O edital pede “fundamentos \textbf{e aplicações}” --- ou seja, a mesma
pergunta técnica pode vir embrulhada num cenário concreto (um órgão
público, um banco, uma prefeitura). Reconhecer o padrão por trás do
cenário resolve a questão:

\begin{itemize}
\item \textbf{Triagem/priorização de atendimento} (fila de posto de
saúde, concessão de benefício, análise de currículo): tipicamente
\textbf{classificação supervisionada} (existe um rótulo conhecido:
aprovado/negado, prioritário/não) --- e o cenário costuma testar se há
\textbf{direito a explicação/revisão humana} da decisão.
\item \textbf{Detecção de fraude ou anomalia} (transação suspeita, glosa
irregular): tipicamente \textbf{não supervisionado} ou semissupervisionado
--- porque casos rotulados de fraude são raros e o padrão fraudulento
\textbf{muda com o tempo} (concept drift, ver Capítulo
\ref{cap:orfaos-teoria}).
\item \textbf{Chatbot/atendimento ao cidadão via LLM}: testa se o
candidato reconhece o risco de \textbf{alucinação} e a necessidade de
\textbf{validação humana} antes que a resposta do modelo produza efeito
sobre um direito do cidadão.
\item \textbf{Score de crédito/seguro/benefício por algoritmo}: cenário
clássico de \textbf{viés algorítmico} --- se o histórico de treino
reflete desigualdade passada, o modelo tende a reproduzi-la, mesmo sem
usar diretamente o atributo sensível (via proxy).
\end{itemize}
\end{conceito}

\section{Interseção IA $\times$ ESG: energia e sustentabilidade}

\begin{conceito}[Por que IA e sustentabilidade aparecem juntas em prova]
O crescimento do uso de IA em escala tem um custo energético real, e é
exatamente o tipo de interseção que uma prova de “Atualidades e IA” gosta
de explorar --- une os dois temas mais recorrentes do bloco:
\begin{itemize}
\item O consumo de um data center se divide, a grosso modo, entre
\textbf{processamento} (a maior fatia) e \textbf{resfriamento}
(refrigerar o hardware para não superaquecer é, depois do próprio
processamento, o maior custo energético de operar um data center em
escala).
\item \textbf{PUE} (\textit{Power Usage Effectiveness}): métrica-padrão de
eficiência energética de data center --- razão entre a energia total
consumida e a energia usada especificamente pela TI. Quanto mais próximo
de 1,0, mais eficiente (menos energia “desperdiçada” em overhead como
refrigeração).
\item Como parte da energia elétrica que alimenta data centers ainda vem
de fontes fósseis, a expansão acelerada de IA \textbf{pressiona} as metas
de descarbonização --- é a mesma preocupação de mudança climática que já
aparece na Parte 1 deste capítulo, agora aplicada especificamente à
infraestrutura digital.
\end{itemize}
\end{conceito}

\begin{cuidado}[Não caia no absoluto em nenhuma das duas direções]
Dizer que a infraestrutura de IA usa “majoritariamente energia limpa” é
otimismo não sustentado (a menos que o texto-base afirme isso
explicitamente); dizer que IA é “incompatível com sustentabilidade” é o
absoluto oposto, igualmente errado. O tratamento realista do tema é o
\textbf{dilema}: benefício tecnológico \textbf{e} custo ambiental
coexistindo, não um veredito fechado num sentido ou no outro.
\end{cuidado}

\section{Roteiro de revisão do capítulo}

\begin{regrapratica}[Os pares que este capítulo mais inverte]
Supervisionado (rotulado) $\times$ não supervisionado (sem rótulo); viés
alto/underfitting (modelo simples demais, erra em tudo) $\times$ variância
alta/overfitting (modelo decorou o treino); precisão (foco no falso
positivo) $\times$ recall (foco no falso negativo); F1 é média
\textbf{harmônica}, não aritmética; AI Act (\textbf{em vigor}) $\times$ PL
2338 (\textbf{ainda projeto}); explicabilidade $\times$ acurácia (eixos
independentes); “não usa o atributo sensível” não é o mesmo que “livre de
viés” (cuidado com proxy).
\end{regrapratica}
